\documentclass[10pt]{article}
\usepackage{template}
\usepackage{amsmath,amssymb,microtype,tikz,graphicx,xurl}
\usetikzlibrary{arrows.meta,positioning,calc}
\setlength{\emergencystretch}{3em}
\raggedbottom
\SetKeys{date=2026-11-01,vol=019,title={每日有趣新闻｜边界、证据与未来},author={ChatGPT 编辑部},note={\small 采编截止：北京时间11月1日上午；主要追踪10月30—31日发布或发生的事件。政府声明、媒体报道、作者证明主张与本刊分析分别归属；新闻图片与数据图不得混作现场证据。}}
\renewcommand{\newsitem}[3]{\par\medskip{\large\bfseries\color{newsblue}#1\par}{\footnotesize\color{newsgray}#2\par}\smallskip\color{black}#3\par\medskip{\color{newsgray}\rule{\linewidth}{0.25pt}}\par}
\newcommand{\newslead}[3]{\par\medskip{\LARGE\bfseries\color{newsblue}#1\par}\vspace{1mm}{\small\color{newsgray}#2\par}\smallskip\color{black}#3\par\medskip{\color{newsgray}\rule{\linewidth}{0.45pt}}\par}
\newcommand{\newssub}[1]{\par\smallskip\noindent{\bfseries\color{newsblue}#1}\quad}
\begin{document}
\begin{documents}
\raggedcolumns
\newslead{西撒哈拉的联合国任务到期：安理会为什么在两份决议之间都说了「不」？}{国际政治深度｜10月30日安理会表决；10月31日任务期限届满；联合国决议、美联社、路透社与冲突双方声明交叉核查}{
西撒哈拉的联合国特派团MINURSO，正式名称是「联合国西撒哈拉全民投票特派团」，自1991年起在这片前西班牙殖民地工作。它原本承载两项任务：观察停火，并协助为领土未来举行全民投票。35年过去，投票从未实现。10月30日，安理会对美国与俄罗斯分别提出的续期方案都未能达成决议；依据2025年第2797号决议，原有授权截止10月31日。进入11月1日时，本刊核查到的公开材料尚无新决议，因此\textbf{可确认的是任务授权未获续期，不能进一步断言所有人员已撤离、联合国已经放弃该议题}。\cite{UN2797,APWestern,ReutersWestern}

\newssub{两份草案为什么不能互相替代？}美国提出的不只是机械延长一年，而是改变任务性质：把摩洛哥2007年提出的自治方案确立为谈判的唯一可行基础，逐步削减军事观察任务，并考虑改为政治特派团。该方案把西撒哈拉置于摩洛哥主权之下，同时允许一定程度的地方自治。俄罗斯则主张不改变原有使命，先延长MINURSO一年，再继续政治磋商。两种文本的争议不只是行政名称，而是\textbf{是否在正式谈判之前就排除了独立这一结果}。\cite{APWestern,ReutersWestern}

\begin{center}\resizebox{\columnwidth}{!}{\begin{tikzpicture}[x=0.33cm,y=0.42cm, font=\scriptsize]
\node[anchor=west,font=\bfseries] at (0,5.6) {安理会10月30日两项草案表决（15国）};
\foreach \y/\name/\a/\b/\c in {3.6/美国草案/11/1/3,1.6/俄罗斯草案/3/1/11}{
 \node[anchor=east] at (7.5,\y+0.5) {\name};
 \fill[newsblue!75] (8,\y) rectangle ({8+\a},\y+1);
 \fill[red!65] ({8+\a},\y) rectangle ({8+\a+\b},\y+1);
 \fill[newsgray!40] ({8+\a+\b},\y) rectangle (23,\y+1);
 \node[white] at ({8+\a/2},\y+0.5) {\a};
 \node[anchor=west] at (24,\y+0.5) {赞成／反对／弃权};
}
\fill[newsblue!75] (0,0) rectangle (1,0.5);\node[anchor=west] at (1.4,0.25) {赞成};
\fill[red!65] (7,0) rectangle (8,0.5);\node[anchor=west] at (8.4,0.25) {反对};
\fill[newsgray!40] (14,0) rectangle (15,0.5);\node[anchor=west] at (15.4,0.25) {弃权};
\end{tikzpicture}}\par\vspace{1mm}{\scriptsize\color{newsgray}图1｜编辑部按美联社及路透社报道绘制；美国草案遭俄否决，俄草案仅3票赞成，未达9票门槛。非现场图。}\end{center}\color{black}

10月30日美国草案获得11票赞成、1票反对、3票弃权，却因俄罗斯作为常任理事国投反对票而遭否决。俄罗斯草案随后仅获3票赞成、1票反对、11票弃权，没有达到安理会15国中至少9票赞成的通过门槛。这里存在容易混淆的程序细节：\textbf{美国确实投了反对票，但俄罗斯草案即使没有美国否决也不会通过}；称两份草案「各遭一次否决」会掩盖第二份文本支持票不足这一独立原因。\cite{APWestern,ReutersWestern}

\newssub{领土、自治与自决。}1975年西班牙撤出后，摩洛哥控制西撒哈拉大部分地区，波利萨里奥阵线主张撒哈拉威人民自决与独立，阿尔及利亚支持该阵线。1991年联合国促成停火，原拟以投票解决地位争议，但谁有资格参加投票长期无法谈妥。2020年以来当地又发生低强度武装冲突。支持摩洛哥方案的人强调自治可能比长期僵局更有现实可行性；反对者则指出，若预先规定主权归属，独立选项将不复存在。这里并非两种宣传语之间的文字争吵，而是\textbf{谈判能否产生当事人尚未预定的政治结果}。\cite{APWestern,UN2797}

\newssub{各方声明应如何读？}摩洛哥王室办公室10月30日发表声明，欢迎MINURSO未获续期，并主张其人员、车辆和设备应在未来几天撤离；还把此事解释为联合国应结束相关议程的信号。这是摩洛哥政府的\emph{政治诉求}，不是联合国正式宣布关闭西撒哈拉档案。波利萨里奥阵线10月31日声明则指责美国方案损害自决权，表示仍愿与联合国合作。两份原始声明具有鲜明的利益立场，本刊不把其中对主权合法性的定性当作已裁定事实。\cite{MoroccoStatement,PolisarioStatement}

\newssub{任务终止之后，危险在哪里？}即便军事观察团没有阻止全部冲突，其巡逻、联络和情况报告仍是外界观察当地安全态势的一条渠道。缺少授权会引出人员安全、设备撤出、停火线监测、难民营人道援助和未来谈判主持权等一系列问题。摩洛哥要求撤离不等于撤离已经完成；波利萨里奥表示继续谈判也不等于双方恢复正式谈判。更值得关注的是联合国是否通过新的技术性安排维持最低限度观察，是否公布撤离计划，以及是否存在独立可核验的平民安全报告。\cite{APWestern,ReutersWestern}

\newssub{编辑观察。}安理会常被批评因常任理事国否决权而陷入瘫痪，但这次事件还展示了另一层制度困境：即使不计否决权，也无法自动得到九票支持。美国希望以明确结果推动谈判，俄罗斯希望维持原框架，两种策略都有其政治目标，却都未能形成足够广泛的共识。真正被悬置的，不仅是一项联合国机构的预算与名义，更是当地居民的安全和未来政治身份问题。\cite{APWestern,UN2797}
}

\newslead{一份数学证明到底该由谁判对？「拉马努金挑战」把AI评卷人也送上考场}{数学与AI深度｜Sorani等人10月29日arXiv:2610.40551；作者自报实验结果，尚未同行评议}{
让大语言模型写一份数学证明，难点不只在于它能不能写出来，更在于\textbf{谁来判定这份证明是对的}。如果把另一台大语言模型放在评卷席上，可能出现一套流畅但错误的论证，被同样流畅的评语批准。10月29日，Alan Sorani、Ashvni Narayanan、Michael Shalyt等人在arXiv提交论文，报告「拉马努金挑战」（Ramanujan Challenge）的参赛结果及一条结合自然语言、计算机代数系统（CAS）与Lean形式化的评估流程。它研究的是研究级数论题，不是中学奥数的选择题。\cite{Ramanujan}

\newssub{定义：答案能验，证明未必能验。}挑战收录十道数论问题，其中八道已有组织者掌握但尚未发表的证明，另两道是开放猜想。题目要求给出某些数学常数的具体公式。你可以用高精度数值计算比较一个公式在前几百位小数是否正确，但有限精度的一致并不能证明恒等式对所有情况成立。最朴素的反例是两个表达式在大量样本点相同，却在未检查的点上不同。因此「数值可检验」与「证明可验证」是两层不同的证据。\cite{Ramanujan}

\newssub{例子：不同工具究竟负责什么？}自然语言证明能够表达研究者的直觉和结构性论证，但可能遗漏假设、把近似当等式或使用尚未证明的引理。CAS可以精确操作符号表达式、复算有限求和、验证代数恒等式，却必须确认输入命题与实际研究命题一致。Lean等交互式定理证明器能够让一个小型内核检查形式化证明的每一步推理，但前提是形式化陈述确实忠实对应原问题，并且实际完成了内核检查。工具的多样性增加了发现错误的机会，却不等于三台工具投票即可自动获得数学真理。\cite{Ramanujan}

\begin{center}\begin{tikzpicture}[>=Latex,node distance=1.2mm, every node/.style={font=\scriptsize},box/.style={draw=newsblue,rounded corners,align=center,minimum width=4.2cm,minimum height=0.55cm,fill=newsblue!5}]
\node[box] (a) {研究级数论问题与候选证明};
\node[box,below=of a] (b) {自然语言证明／CAS代码／Lean文件};
\node[box,below=of b] (c) {语义检查、计算复算、形式化内核核验};
\node[box,below=of c] (d) {专家复核与公开可复现材料};
\draw[->,thick] (a)--(b);\draw[->,thick] (b)--(c);\draw[->,thick] (c)--(d);
\end{tikzpicture}\par\vspace{1mm}{\scriptsize\color{newsgray}图3｜编辑部概念流程，非论文原始算法图；不同证据链的检查力度不相同，Lean文件须实际通过内核才算机器核验。}\end{center}\color{black}

\newssub{困难：评审量本身爆炸。}作者统计收到228份提交，共计4802页自然语言证明与约107.9万行Lean、Python、Mathematica、Sage等代码。面对如此多材料，逐页人工审读需要大量专业时间，而仅用单次模型回答「正确或错误」又有很高的误判风险。论文声称，新的分层评估流程把自然语言证明的\textbf{假阳性率由23\%降至6\%}，辅助代码判断准确率从46\%提高至58\%。这些数字属于作者在其数据与评估条件下的结果；它们不是所有数学证明场景中的通用准确率，58\%也远非可以代替专家复核的程度。\cite{Ramanujan}

\newssub{方法：把一份证明拆成可审计的证据链。}评估流程尝试分别处理数学语义、计算支撑和形式化文件。对自然语言论证，检查关键引理是否被真正使用、是否偷偷强化假设；对代码，复算结果、检查数值精度和边界条件；对Lean，检查陈述、依赖及内核是否接受。作者还报告部分正确投稿的思路与官方解法不同，并声称至少一项已公开开放猜想获得证明。此处必须严格区分：\textbf{「作者认定参赛证明正确」不是本刊已独立审阅其证明}，而开放问题被解决的学界共识仍需要可公开复现材料与进一步检查。\cite{Ramanujan}

\newssub{为什么这比模型排行榜有意思？}研究级数学的一个特点是，正确的结果可能非常简短，而证据链非常漫长。若用单一「模型得分」衡量，会把出题、证明、代码与审稿四种能力混在一起。作者观察到，成功投稿者包括没有高等数学学历的参与者，也有年轻独立研究者；这是值得研究的知识门槛变化，但尚不能推出「AI已能独立解决任何开放问题」。真正可能改变科研实践的是可审计的协作：机器帮助探索与形式化，人类负责问题设定、解释和优先级判断，最后由明确的验证机制决定哪些结论可以进入知识体系。\cite{Ramanujan}

\newssub{验证状态。}这是一篇2026年10月29日发布的arXiv预印本。论文中的竞赛统计、正确性判断和流程效果均为作者报告；本刊未取得独立复现实验或同行评议完成的证据，亦未逐行编译论文所涉Lean证明。它提供了值得认真测试的评估方法，而不是对AI数学能力的终审判决。\cite{Ramanujan}
}

\newsitem{制造业PMI降至49.3：扩散指数的「收缩」究竟意味着什么？}{中国经济重点｜国家统计局10月31日公布原始数据及官方解读}{
10月31日，国家统计局公布制造业采购经理指数（PMI）为49.3，比9月下降0.8个百分点；非制造业商务活动指数与综合PMI产出指数均为49.7。9月制造业PMI为50.1，本月重新跌至50的分界线以下。PMI不是当月产量的同比增长率，而是企业对订单、生产、库存、雇员等情况的调查形成的\textbf{扩散指数}。49.3并不等于工业总产出下降0.7\%，更不能直接理解为GDP下降0.7\%。\cite{NBSOct,NBSExplain}

\begin{center}\begin{tikzpicture}[x=0.60cm,y=0.24cm,font=\scriptsize]
\node[anchor=west,font=\bfseries] at (0,15.2) {10月PMI及行业分化（\%）};
\draw[newsgray,dashed] (8,0) -- (8,14);
\node[above] at (8,14) {50 荣枯线};
\foreach \y/\lab/\v in {12/制造业/49.3,10/非制造业/49.7,8/综合产出/49.7,6/高技术制造业/51.3,4/大型制造企业/50.2,2/小型制造企业/47.5}{
\node[anchor=east] at (4.2,\y+0.6) {\lab};
\pgfmathsetmacro{\w}{(\v-45)*1.6}
\fill[newsblue!65] (0,\y) rectangle (\w,\y+1.2);
\node[anchor=west] at ({\w+0.2},\y+0.6) {\v};
}
\end{tikzpicture}\par\vspace{1mm}{\scriptsize\color{newsgray}图2｜2026年10月扩散指数，横向条形以45为截断基线（非零起点），仅为显示50附近差异；来源：国家统计局10月31日。}\end{center}\color{black}

\newssub{生产与需求同时回落。}生产指数49.8，较上月下降1.9个百分点；新订单49.4，下降1.1个百分点。原材料库存48.4，从业人员48.4，供应商配送时间50.2。多数构成指数处于50以下，意味着相较上月报告改善的企业比例不足以抵消报告恶化的比例，但不是所有企业都在减产。官方解读提到国庆假期前需求提前释放与复杂国际环境；这是一种可能的解释，而非经独立因果识别确认的唯一原因。\cite{NBSOct,NBSExplain}

\newssub{企业规模分化值得注意。}大型企业50.2，仍在扩张区间；中型企业48.4，小型企业47.5，且后两者较上月分别下降1.3和1.4个百分点。高技术制造业51.3、装备制造业50.6，表现相对较好；消费品行业48.8、高耗能行业47.3则偏弱。不能把「高技术行业仍扩张」直接推广为整个制造业繁荣，也不能因为总指数低于50便忽略内部结构的差异。\cite{NBSExplain}

\newssub{下一步看什么？}应结合工业增加值、企业利润、出口订单、社会融资、居民消费和就业数据观察PMI下行是否持续。样本调查能比部分实物统计更快反映方向变化，却可能受到季节、假期和企业主观判断影响。与其围绕49.3这个单一数字争论「衰退还是复苏」，不如追踪新订单是否重新站上50，以及中小企业能否同步改善。\cite{NBSOct}
}

\newsitem{在地球上「住火星」378天：NASA封闭舱实验究竟测什么？}{航天与人因工程重点｜NASA10月27日活动预告，10月31日预定结束；结束报道与长期实验数据仍待核查}{
NASA的第二次CHAPEA（Crew Health and Performance Exploration Analog）模拟火星任务，安排四名志愿者从2025年10月19日起在休斯敦约1700平方英尺的3D打印栖居设施中生活378天。NASA原定于10月31日美国东部时间15时举行出舱活动，换算北京时间为11月1日凌晨3时。已有媒体报道志愿者按计划离舱；但本刊在截稿时尚未取得NASA发布的完整任务后评估报告，故不把活动举行等同于科学数据已完成分析。\cite{NASACHAPEA,NASAOneYear}

\begin{center}\begin{tikzpicture}[>=Latex,font=\scriptsize]
\draw[thick,rounded corners,fill=newsblue!5] (0,0) rectangle (7.2,2.1);
\node[font=\bfseries] at (3.6,1.83) {地面封闭栖居舱（示意）};
\foreach \x/\t in {0.9/医疗,2.65/种植,4.4/维护,6.15/运动}{\draw[fill=white] (\x-0.7,0.55) rectangle (\x+0.7,1.35);\node at (\x,0.95) {\t};}
\draw[<->,thick] (3.6,-0.12)--(3.6,-0.8);\node[anchor=north] at (3.6,-0.8) {延迟通信、限量资源、长期隔离};
\end{tikzpicture}\par\vspace{1mm}{\scriptsize\color{newsgray}图4｜NASA CHAPEA研究要素编辑示意，非真实舱室平面图，不按比例；NASA公布舱室面积约1700平方英尺。}\end{center}\color{black}

\newssub{模拟的不是火星重力。}四名成员Ross Elder、Ellen Ellis、Matthew Montgomery和James Spicer分别承担指挥、医疗、科学与工程职责。他们进行模拟舱外活动、种植作物、设施维护、锻炼和日常医疗任务。实验刻意引入隔离、通信延迟和资源限制，以观察团队如何分配工作时间、处理意外、保持心理健康与执行协作程序。真正的火星还有低重力、宇宙辐射、复杂尘埃、无法快速撤离等问题，休斯敦封闭舱不能完全重现。\cite{NASACHAPEA}

\newssub{为什么要把人关这么久？}短期实验能测试一个设备或流程能否使用，长期隔离才容易暴露重复劳动、单调生活、睡眠节律、成员关系与任务效率之间的反馈。例如，食物种植系统不仅要计算产量，还要计算维护工时和失败时的替代供给；延迟通信意味着地面控制中心不能像近地轨道任务那样实时介入。NASA称研究结果将帮助设计未来月球和深空任务中的食物系统、医疗支持、通信带宽与人员工作安排。\cite{NASACHAPEA,NASAOneYear}

\newssub{证据的边界。}四人小样本不是可直接外推到所有宇航员的统计总体；志愿者知道自己仍在地球上，心理预期也不同于真正远离地球的航行。该任务最有价值的是与其他地面模拟、国际空间站实测及未来任务形成交叉验证，而不是把「378天成功生活」当作人类已掌握火星移民技术。\cite{NASACHAPEA}
}

\newsitem{消费基础设施不是又建一批商场：五部门新意见如何理解？}{中国公共政策重点｜国家发展改革委等部门10月30日原始政策文件}{
国家发展改革委等部门10月30日发布《关于加强消费基础设施建设促进消费投资良性互动的指导意见》。文件涉及农贸市场、乡镇商贸中心、商业综合体、步行街、文化体育设施及相关设备更新。它试图通过改善消费场所的交通、仓储、卫生、照明、安全与服务条件，降低消费者使用成本，同时形成建设投资。文件的政策目标是「消费与投资良性互动」，不是公布一笔已经支出的专项资金，也不是要求所有地方无差别新建大型商业地产。\cite{NDRCConsume}

\newssub{公共服务与商业回报并不总一致。}老旧农贸市场的排水、冷链与消防改造可能提高食品安全和经营效率，但收费能力有限；商业街区引入文旅体业态可能吸引客流，却也有同质化和短期网红项目的风险。地方政府需要判断人口分布、消费能力、现有设施空置率与长期维护成本。一个场馆的开工投资属于当期需求，能否在未来持续创造服务价值，则需要看实际使用频率与经营现金流。\cite{NDRCConsume}

\newssub{应当如何验收？}比起统计新增商场数量，更有意义的指标是居民平均到达时间、无障碍覆盖、食品冷链损耗、消防事故率、老旧设施使用率及小商户负担。政策文本提供了建设方向，尚未提供这些结果指标的独立事后评价。地方落实时，防止重复建设、保护既有小商户和公开招投标信息，将直接决定它究竟是改善消费环境还是制造新的闲置资产。\cite{NDRCConsume}
}

\newsitem{信托资管新规：为什么「卖者尽责、买者自负」必须连着读？}{金融监管重点｜10月30日金融监管总局答记者问；2027年1月1日起施行}{
国家金融监督管理总局公布《资产管理信托管理办法》，将资产管理信托明确定位为私募资产管理产品，要求面向合格投资者非公开发行，投资者数量不超过200人。它强调「卖者尽责、买者自负；卖者失责、依法赔偿」。后半句尤其重要：投资者应承担正常市场风险，不意味着销售机构可以隐瞒风险、虚假宣传或承诺保本。\cite{NFRAtrust}

监管规定禁止任何形式的保本保收益承诺，并限制代销渠道；对较高风险产品要求更严格的投资者资格和最低投资金额。除特定例外外，单一产品投资同一资产或同一资管产品的资金不得超过产品净资产的25\%。还要求不同信托产品单独管理、建账与核算，禁止以不透明资金池掩盖基础资产真实收益及风险。\cite{NFRAtrust}

这些规则的共同目标是减少通道业务、期限错配与关联交易中的信息不对称。新规并不能消灭信用违约，反而可能使部分长期被平滑的净值波动更充分暴露。存量业务按照「新老划断」等安排逐步整改，监管部门表示将避免对实体经济造成过大冲击。具体效果要看穿透披露、违约处置、投资者赔偿和监管执行，而非只看条文是否严厉。\cite{NFRAtrust}
}

\newsitem{欧元区第三季度增长0.4\%：初值比「欧洲经济好坏」更具体}{欧洲经济重点｜Eurostat 10月30日原始统计公报}{
欧盟统计局10月30日发布2026年第三季度GDP初步估算：经季节调整，欧元区环比增长0.4\%，欧盟整体环比增长0.5\%；第二季度分别为0.6\%和0.7\%。同比则分别增长1.3\%和1.5\%。环比增速放缓，并不意味着GDP水平环比下降；同比增长也不能与环比增长直接相加。\cite{EurostatGDP}

\begin{center}\begin{tikzpicture}[x=0.8cm,y=0.65cm,font=\scriptsize]
\node[anchor=west,font=\bfseries] at (0,3.7) {第三季度GDP环比增长初值（\%）};
\foreach \y/\lab/\v in {2.6/欧元区/0.4,1.2/欧盟/0.5}{
\node[anchor=east] at (1.4,\y+0.25) {\lab};
\fill[newsblue!65] (1.7,\y) rectangle ({1.7+\v*8},\y+0.55);
\node[anchor=west] at ({1.9+\v*8},\y+0.25) {\v\%};}
\end{tikzpicture}\par\vspace{1mm}{\scriptsize\color{newsgray}图5｜欧盟统计局10月30日初步估算；两组均为经季节调整环比，非同比；基线为0。}\end{center}\color{black}

该机构强调，快速估算使用的数据尚不完整，未来可能修订。对欧盟经济的判断还需观察成员国之间的差异、就业、工资和能源价格，而不能只把21个欧元区成员的总量平均值视为每个国家的经历。与中国10月PMI比较时更须谨慎：一个是季度国民账户的实际产出增长率，一个是月度企业调查的扩散指数，不能因为数字大小接近就认为它们反映同一种变化。\cite{EurostatGDP}
}

\newsitem{黑海局部停火又被提出：保护粮船不等于结束俄乌战争}{外交简讯｜路透社10月31日亚特兰大G20外长会现场报道}{
美国国务卿鲁比奥10月31日表示，美国正在研究土耳其与印度提出的黑海局部停火建议，以减少对粮食和能源运输船舶的袭击。加拿大外长阿南德也称，已与乌克兰、印度等讨论海运安全及冬季电网设施遭袭问题，但俄罗斯尚未显示愿意推进的态度。G20外长会议未能形成联合声明。\cite{ReutersBlackSea}

这里的「兴趣增加」仍是外交官对谈判可能性的描述，不是俄乌签署停火文本。即使未来达成海运协议，也需明确适用船只、港口、监测机构与违约处理；黑海航运安全与陆上战事停止并非同一回事。\cite{ReutersBlackSea}
}

\newsitem{数学拾趣｜六维、七维的「挡住视线」问题}{数论与离散几何简讯｜Vineet Kahlon；arXiv:2610.37711，10月26日提交、30日新列表；未同行评议}{
在$n$维空间中，把同半径闭球放在格点$(1/2,\ldots,1/2)+\mathbb Z^n$上。问题是球至少要多大，才能让从原点射向正正交象限的每一条射线都撞到某个球？这称为球的\textbf{视线遮挡问题}。低维答案早已知道，高维的极值则牵涉格点几何与精确计算。\cite{Kahlon}

Kahlon在预印本中声称得到六维和七维的精确极值$D_6=107/312$、$D_7=73/180$，并给出相应的极值方向。作者的方法把无限优化问题归约为有限搜索，利用对偶格、Banaszczyk转移定理与精确有理数计算形成搜索证书。\textbf{这些数是论文所定义的平方距离型极值$D_n$，不是球半径本身}。本刊尚未复跑全部计算证书或核查证明细节；预印本也尚无可确认的同行评议完成记录。\cite{Kahlon}
}


\newsitem{哈勃还没退休：NASA重新征集商业维修方案，但不是已经决定派飞船}{航天工程重点｜NASA10月30日公告；10月29日发布信息征询（RFI）}{
NASA于10月30日宣布，已在29日发布面向产业界的信息征询，探索是否可以为运行逾35年的哈勃空间望远镜提供商业在轨服务。候选方向包括提高轨道高度、维修或提升部分能力，以及在寿命终点进行受控再入或轨道处置。NASA估计，如果不采取新措施，哈勃可能在2033—2034年左右重返大气层；如果要实施服务任务，时间窗口可能需要落在2032年之前。\textbf{这是一份征询能力与方案的文件，不是已批准预算、选定承包商或确定发射日期。}\cite{NASAHubble}

\newssub{为什么升高轨道有意义？}低地球轨道并非完全真空，稀薄大气对卫星产生阻力，长期使轨道高度降低。哈勃没有自主维持轨道的推进系统，因此即使科学仪器继续工作，轨道寿命仍会缩短。通过外部航天器捕获并提升轨道，有机会增加剩余运行年限；但接近一台精密、昂贵、并非为任意商业飞船对接而设计的望远镜，本身可能引入碰撞、姿态失控和结构损伤风险。\cite{NASAHubble}

\newssub{三种目标并不等价。}「升轨」主要解决阻力和寿命；「维修升级」可能处理设备故障或增加观测能力；「受控处置」则是减少未来无控制再入的地面风险。它们所需机械接口、导航精度、保险和失败预案都不同。NASA希望在11月30日之前收到产业界回复，随后才能评估技术成熟度与成本收益。此前哈勃已接受多次航天飞机维修，但商业机器人服务的操作环境不同。\cite{NASAHubble}

\newssub{怎样评价是否值得做？}比较方案不能只看延寿几年，还应考虑剩余科学产出、观测独特性、对下一代旗舰望远镜的资源挤出效应，以及失败时是否反而加速任务终结。另一边，哈勃与韦布在波段和观测方式上并非简单替代关系。眼下最重要的信息增量是NASA正式征询方案；「哈勃获救」还远远谈不上。\cite{NASAHubble}
}

\par\medskip{\large\bfseries\color{newsblue}来源、图注与编辑说明}\par
{\scriptsize 本期主动检索国际时政外交、战争、国内政策、社会议题、科技与AI、Rust/C++、Linux、物理、化学、航天、文化和arXiv数学（数论、代数、几何、组合、逻辑、概率、Lean/mathlib）。未找到足够可靠的当日增量者不强行刊登。五张图均为编辑部以可追溯公开数据或概念定义绘制的TikZ矢量图，实际源码随本期归档；未使用AI合成新闻现场图。NASA CHAPEA的真实现场照片可在NASA原始发布页核查，但本期不以未获得独立文件的网页缩略图冒充归档原图。所有媒体来源和作者主张均按其原始归属引用；10月31日发生的事件不等于11月1日已完成全部独立核验。}\par
\begingroup\fontsize{6}{7.2}\selectfont
\renewcommand{\refname}{参考文献}
\bibliographystyle{plainurl}\bibliography{refs}
\endgroup
\end{documents}
\end{document}
